\documentclass[10pt,a4paper]{amsart}
\usepackage[margin=19mm]{geometry}
\usepackage{amsmath,amssymb,mathtools}
\usepackage[scr=rsfs,cal=euler]{mathalfa}
\usepackage{xfrac}
\usepackage[unicode,hidelinks,bookmarks=false]{hyperref}
\newcommand{\ie}{i.e.\ }
\newcommand{\cf}{cf.\ }
\newcommand{\resp}{respectively\ }
\newcommand{\Subsection}{\subsection}
\providecommand{\nobreakdash}{\nobreak\hbox{-}}
\setlength{\emergencystretch}{2em}
\multlinegap0pt
\usepackage{amssymb}
\usepackage{esint}
\usepackage{mathtools}
\usepackage{xcolor}
\newtheorem{proposition}{Proposition}
\newcommand{\N}{\mathbb{N}}
\newcommand{\R}{\mathbb{R}}
\title{Asymptotic minimality of one-dimensional transition profiles in Aviles-Giga type models: an approach via $1$-currents}
\author{Radu Ignat, Roger Moser}
\date{Source: arXiv:2508.13753v2 (CC BY 4.0)}
\begin{document}
\maketitle

\noindent\emph{Source and licence.} Asymptotic minimality of one-dimensional transition profiles in Aviles-Giga type models: an approach via $1$-currents. Version arXiv:2508.13753v2 (CC BY 4.0), distributed by arXiv under the Creative Commons Attribution 4.0 International licence, http://creativecommons.org/licenses/by/4.0/, as recorded in the arXiv licence metadata for this exact version and shown on the abstract page https://arxiv.org/abs/2508.13753v2. Full licence notice: \texttt{licenses/arxiv-2508.13753-CC-BY-4.0.txt}.\par\smallskip

\noindent\textit{Editorial scope.} Counted unit: the article's proposition prp:polynomial (source0.tex lines 3347--3359). The article's complete original proof of it is delivered with it (source0.tex lines 3361--3483, one \texttt{proof} environment of 123 lines). No mathematics is written, repaired or re-derived: the statement and the proof are the article's own lines, in the article's own notation, and no retained line is substituted or re-flowed.  No further source-local context is needed: every cross-reference of every retained line resolves inside the two delivered spans.  Nothing was omitted from any retained span, so the package carries no omission record.  No retained line contains a citation, so the wrapper carries no bibliography and no bibliography entry.  No retained line contains a figure environment, an image-inclusion command or drawing code, so no image is delivered and the package declares no asset dependency.  Editorial formatting only, in addition: an AMS \texttt{amsart} preamble that restates the article's own declarations and macros used by the retained lines, namely the environments \texttt{proposition} on separate counters, together with the standard packages those lines need.  The retained environments print in this document as Proposition 1 in order of appearance, whereas the article prints the same items with the numbering of its own full document.  The complete native source of the article as distributed in the arXiv e-print for this exact version is bundled unchanged as \texttt{sources/183504/source0.tex}.
\begin{proposition} \label{prp:polynomial}
For any $n \in \N_0$, there exists a polynomial $P \colon \R^2 \to \R$ such
that
\begin{equation} \label{eq:second-derivatives-polynomial}
\left|\frac{\partial^2 P}{\partial y_1^2}(y)\right| \le |y|^{2n} \bigl|1 - |y|^2\bigr| \quad \text{and} \quad \left|\frac{\partial^2 P}{\partial y_2^2}(y)\right| \le |y|^{2n} \bigl|1 - |y|^2\bigr|
\end{equation}
for every $y \in \R^2$ and
and
\begin{equation} \label{eq:equality-polynomial}
\frac{\partial^2 P}{\partial y_2^2}(0, y_2) = y_2^{2n} (1 - y_2^2)
\end{equation}
for every $y_2 \in \R$.
\end{proposition}

\begin{proof}
For $n = 0$, a suitable polynomial is
\[
P(y) = \frac{|y|^2}{2} - \frac{|y|^4}{12} - \frac{y_1^2 y_2^2}{3}.
\]

We now assume that $n \ge 1$. We look for a polynomial such that
\[
\frac{\partial^2 P}{\partial y_2^2}(y) = (1 - |y|^2) \left(y_2^{2n} + \sum_{k = 0}^{n - 1} c_k y_1^{2n - 2k} y_2^{2k}\right)
\]
for certain coefficients $c_0, \dotsc, c_{n - 1}$. Setting $c_n = 1$, we
can write
\[
\begin{split}
\frac{\partial^2 P}{\partial y_2^2}(y) & = (1 - |y|^2) \sum_{k = 0}^n c_k y_1^{2n - 2k} y_2^{2k} \\
& = (1 - y_1^2) \sum_{k = 0}^n c_k y_1^{2n - 2k} y_2^{2k} - \sum_{k = 0}^n c_k y_1^{2n - 2k} y_2^{2k + 2}.
\end{split}
\]
A possible solution is
\[
\begin{split}
P(y) & = (1 - y_1^2) \sum_{k = 0}^n \frac{c_k}{(2k + 2)(2k + 1)} y_1^{2n - 2k} y_2^{2k + 2} \\
& \quad - \sum_{k = 0}^n \frac{c_k}{(2k + 4)(2k + 3)} y_1^{2n - 2k} y_2^{2k + 4} \\
& \quad + \frac{y_1^{2n + 2}}{(2n + 2)(2n + 1)} - \frac{y_1^{2n + 4}}{(2n + 4)(2n + 3)} \\
& = \sum_{k = 0}^{n - 1} \frac{c_k}{(2k + 2)(2k + 1)} y_1^{2n - 2k} y_2^{2k + 2} + \frac{y_1^{2n + 2} + y_2^{2n + 2}}{(2n + 2)(2n + 1)} \\
& \quad - \sum_{k = 1}^{n - 1} \frac{c_k + c_{k - 1}}{(2k + 2)(2k + 1)} y_1^{2n - 2k + 2} y_2^{2k + 2} - \frac{c_0}{2} y_1^{2n + 2} y_2^2 \\
& \quad - \frac{1 + c_{n - 1}}{(2n + 2)(2n + 1)} y_1^2 y_2^{2n + 2} - \frac{y_1^{2n + 4} + y_2^{2n + 4}}{(2n + 4)(2n + 3)}.
\end{split}
\]

We want to impose the symmetry condition $P(y_1, y_2) = P(y_2, y_1)$
(so that the above condition on $\frac{\partial^2 P}{\partial y_2^2}$
automatically gives a similar condition for
$\frac{\partial^2 P}{\partial y_1^2}$). This requires that
\begin{equation} \label{eq:coefficients1}
\frac{c_0}{2} = \frac{1 + c_{n - 1}}{(2n + 2)(2n + 1)}
\end{equation}
and
\begin{equation} \label{eq:coefficients2}
\frac{c_k}{(2k + 2)(2k + 1)} = \frac{c_{n - k - 1}}{(2n - 2k)(2n - 2k - 1)}, \quad k = 0, \dotsc, n - 1,
\end{equation}
and
\begin{equation} \label{eq:coefficients3}
\frac{c_k + c_{k - 1}}{(2k + 2)(2k + 1)} = \frac{c_{n - k} + c_{n - k - 1}}{(2n - 2k + 2)(2n - 2k + 1)}, \quad k = 1, \dotsc, n - 1.
\end{equation}
(It may appear at first that there are too many equations for the $n$ variables
$c_0, \dotsc, c_{n - 1}$, but there is some repetition here.
Once the redundant equations are discarded, it is easy to see that there
is a unique solution.)

The combination of \eqref{eq:coefficients2} and \eqref{eq:coefficients3} gives
\[
\begin{split}
c_k + c_{k - 1} & = \frac{(2k + 2)(2 k + 1)}{(2n - 2k + 2)(2n - 2k + 1)}(c_{n - k} + c_{n - k - 1}) \\
& = \frac{(2n - 2k)(2n - 2k - 1)}{(2n - 2k + 2)(2n - 2k + 1)} c_k + \frac{(2k + 2)(2 k + 1)}{2k(2k - 1)} c_{k - 1}
\end{split}
\]
for $k = 1, \dotsc, n - 1$. Thus
\[
\left(1 - \frac{(2n - 2k)(2n - 2k - 1)}{(2n - 2k + 2)(2n - 2k + 1)}\right) c_k = \left(\frac{(2k + 2)(2 k + 1)}{2k(2k - 1)} - 1\right) c_{k - 1}.
\]
We compute
\[
1 - \frac{(2n - 2k)(2n - 2k - 1)}{(2n - 2k + 2)(2n - 2k + 1)} = \frac{8(n - k) + 2}{(2n - 2k + 2)(2n - 2k + 1)}
\]
and
\[
\frac{(2k + 2)(2 k + 1)}{2k(2k - 1)} - 1 = \frac{8k + 2}{2k(2k - 1)}.
\]
Therefore, we obtain the equation
\[
\frac{c_k}{c_{k - 1}} = \frac{(2n - 2k + 2)(2n - 2k + 1)(4k + 1)}{2k(2k - 1)(4(n - k) + 1)}
\]
for $k = 1, \dotsc, n - 1$.

Define
\[
b_k = \frac{c_k}{(2k + 1)\binom{n}{k}}, \quad k = 0, \dotsc, n - 1.
\]
Then
\[
\begin{split}
\frac{b_k}{b_{k - 1}} & = \frac{k(2k - 1)}{(n - k + 1)(2k + 1)} \frac{c_k}{c_{k - 1}} \\
& = \frac{(2n - 2k + 1)(4k + 1)}{(2k + 1)(4(n - k) + 1)} \\
& = \frac{8nk - 8k^2 + 2n + 2k + 1}{8nk - 8k^2 + 4n - 2k + 1}.
\end{split}
\]
We note that $b_k/b_{k - 1} \ge 1$ if, and only if, $k \ge n/2$; and
$b_k/b_{k - 1} \le 1$ if, and only if, $k \le n/2$.
Hence $b_k$, as a function of $k$, is first decreasing, may possibly be constant
for one step, and is then increasing. (If $n = 1$, then this is a vacuous statement, as
we have only $b_0$ in this case.)

From \eqref{eq:coefficients1} and \eqref{eq:coefficients2} for $k = 0$, we also
conclude that
\[
\frac{1 + c_{n - 1}}{(2n + 2)(2n + 1)} = \frac{c_{n - 1}}{2n(2n - 1)}.
\]
Solving this equation, we obtain
\[
c_{n - 1} = \frac{2n^2 - n}{4n + 1}.
\]
It then also follows from \eqref{eq:coefficients2} that
\[
c_0 = \frac{c_{n - 1}}{2n^2 - n} = \frac{1}{4n + 1}.
\]
This means that
\[
b_0 = c_0 = \frac{1}{4n + 1} \quad \text{and} \quad b_{n - 1} = \frac{c_{n - 1}}{2n^2 - n} = \frac{1}{4n + 1}.
\]
We conclude that $b_k \le \frac{1}{4n + 1}$ for all $k = 0, \dotsc, n - 1$,
i.e.,
\[
c_k \le \frac{2k + 1}{4n + 1} \binom{n}{k}.
\]
Because $c_k/c_{k - 1} > 0$ for every $k = 1, \dotsc, n - 1$, it is clear that
$c_k > 0$ for every $k = 1, \dotsc, n$. Therefore,
\[
0 \le \sum_{k = 0}^n c_k y_1^{2n - 2k} y_2^{2k} \le \sum_{k = 0}^n \binom{n}{k} y_1^{2n - 2k} y_2^{2k} = |y|^{2n}.
\]
The inequalities in \eqref{eq:second-derivatives-polynomial} follow.
We also have identity \eqref{eq:equality-polynomial} by construction.
\end{proof}

\end{document}
